% ===== PRÉAMBULE CANONIQUE — à coller VERBATIM en tête de CHAQUE .tex =====
\documentclass[11pt,a4paper]{article}
\usepackage[utf8]{inputenc}\usepackage[T1]{fontenc}\usepackage{lmodern}
\usepackage[french]{babel}
\usepackage{amsmath,amssymb,mathtools}\usepackage{icomma}
\usepackage{siunitx}\sisetup{locale=FR}  % \qty{}{}, \si{}, \ang{}
\usepackage{xcolor}\usepackage{colortbl}
\usepackage{tikz}\usepackage{pgfplots}\pgfplotsset{compat=1.18}
\usepackage[siunitx,europeanresistors]{circuitikz} % circuits (élec.)
\usepackage[most]{tcolorbox}\usepackage{enumitem}\usepackage{array,booktabs}
\usepackage[a4paper,margin=2.2cm]{geometry}
\usetikzlibrary{arrows.meta,calc,angles,quotes,positioning,patterns,%
  backgrounds,shapes.geometric,decorations.pathmorphing,decorations.markings,intersections}
\definecolor{primary}{RGB}{222,49,99}\definecolor{primarydark}{RGB}{150,22,55}
\definecolor{defcol}{RGB}{0,114,178}\definecolor{methcol}{RGB}{52,92,148}
\definecolor{excol}{RGB}{0,135,90}\definecolor{attncol}{RGB}{200,40,55}
\tcbset{boxbase/.style={enhanced,breakable,boxrule=0.9pt,arc=2.5pt,
  left=3mm,right=3mm,top=2.3mm,bottom=2.3mm,fonttitle=\bfseries,coltitle=white}}
% --- boîtes TUTORIEL (noms EXACTS, ne pas renommer) ---
\newtcolorbox{cle}[1][]{boxbase,colback=defcol!6,colframe=defcol,
  title={\textsc{À retenir}\ifx\relax#1\relax\relax\else\textmd{ : \itshape #1}\fi}}
\newtcolorbox{methode}[1][]{boxbase,colback=methcol!7,colframe=methcol,sharp corners,
  title={\textsc{Méthode}\ifx\relax#1\relax\relax\else\textmd{ --- \itshape #1}\fi}}
\newtcolorbox{exemplebox}[1][]{boxbase,colback=excol!5,colframe=excol,
  title={\textsc{Exemple}\ifx\relax#1\relax\relax\else\textmd{ : \itshape #1}\fi}}
\newtcolorbox{piege}[1][]{boxbase,colback=attncol!6,colframe=attncol,
  title={\textsc{Piège}\ifx\relax#1\relax\relax\else\textmd{ : \itshape #1}\fi}}
\newtcolorbox{remarque}[1][]{boxbase,colback=black!4,colframe=black!45,
  title={\textsc{Remarque}\ifx\relax#1\relax\relax\else\textmd{ : \itshape #1}\fi}}
% --- boîtes EXERCICE / CORRIGÉ (noms EXACTS, auto-comptées) ---
\newtcolorbox[auto counter]{exo}[1][]{boxbase,colback=primary!4,colframe=primary,
  title={\textsc{Exercice}~\thetcbcounter\ifx\relax#1\relax\relax\else\textmd{ --- \itshape #1}\fi}}
\newtcolorbox[auto counter]{corrige}[1][]{boxbase,colback=excol!4,colframe=excol,
  title={\textsc{Corrigé}~\thetcbcounter\ifx\relax#1\relax\relax\else\textmd{ --- \itshape #1}\fi}}
\newcommand{\vv}[1]{\vec{#1}}\newcommand{\ds}{\displaystyle}
\newcommand{\R}{\mathbb{R}}\newcommand{\N}{\mathbb{N}}\newcommand{\Z}{\mathbb{Z}}
% (un \usepackage{fancyhdr}+en-tête est possible ; il est ignoré côté web)
% =========================================================================
\begin{document}
\sisetup{per-mode=symbol}

\begin{corrige}[distance de freinage]
\begin{enumerate}[label=\alph*)]
  \item $\qty{108}{\kilo\metre\per\hour}=\dfrac{108}{3{,}6}=\qty{30}{\metre\per\second}$.
  \item Torricelli, $v=0$ : $\Delta x=\dfrac{-v_0^2}{2a}=\dfrac{-30^2}{2\times(-5)}=\dfrac{-900}{-10}=\qty{90}{\metre}$.
  \item $t=\dfrac{-v_0}{a}=\dfrac{-30}{-5}=\qty{6}{\second}$.
\end{enumerate}
\end{corrige}

\begin{corrige}[du repos à la vitesse de croisière]
$v=\qty{72}{\kilo\metre\per\hour}=\qty{20}{\metre\per\second}$.
\begin{enumerate}[label=\alph*)]
  \item Torricelli : $a=\dfrac{v^2-v_0^2}{2\Delta x}=\dfrac{20^2}{2\times100}=\qty{2}{\metre\per\second\squared}$.
  \item $t=\dfrac{v-v_0}{a}=\dfrac{20}{2}=\qty{10}{\second}$.
\end{enumerate}
\end{corrige}

\begin{corrige}[deux phases : accélération puis MRU]
\begin{enumerate}[label=\alph*)]
  \item $t_1=\dfrac{v}{a}=\dfrac{20}{2}=\qty{10}{\second}$ ; distance $x_1=\tfrac12 a t_1^2=\tfrac12\times2\times100=\qty{100}{\metre}$ (ou $\tfrac{v^2}{2a}$).
  \item Phase MRU : $x_2=v\times t=20\times5=\qty{100}{\metre}$.
  \item Distance totale $=x_1+x_2=\qty{200}{\metre}$ (durée totale \qty{15}{\second}).
\end{enumerate}
\end{corrige}

\begin{corrige}[glissade sur un plan incliné]
\begin{enumerate}[label=\alph*)]
  \item $a=g\sin\ang{30}=10\times0{,}5=\qty{5}{\metre\per\second\squared}$.
  \item Torricelli : $v=\sqrt{2a\,\Delta x}=\sqrt{2\times5\times10}=\sqrt{100}=\qty{10}{\metre\per\second}$.
\end{enumerate}
\end{corrige}

\begin{corrige}[lire les paramètres dans l'équation horaire]
On compare à $x=x_0+v_0t+\tfrac12 a t^2$.
\begin{enumerate}[label=\alph*)]
  \item $x_0=\qty{2}{\metre}$ ; $v_0=\qty{4}{\metre\per\second}$ ; $\tfrac12 a=1{,}5 \Rightarrow a=\qty{3}{\metre\per\second\squared}$.
  \item $a>0$ et $v_0>0$ : mouvement accéléré. $v(t)=v_0+at=4+3t$, donc $v(2)=4+6=\qty{10}{\metre\per\second}$.
\end{enumerate}
\end{corrige}

\end{document}
